% !TEX root = ../main.tex
\clearpage
\phantomsection
\section*{8 \textbf{연구의 의의}}
\label{sec:significance-of-research}
\addcontentsline{toc}{section}{8 연구의 의의}

본 연구는 온체인 평판 점수 산출 및 탈중앙화 신용 평가 분야를 여러 의미 있는 방식으로 발전시킨다. 박사 수준에서 의도된 독창적 기여는 승인 기반 지갑 평판을 구별된 유형의 온체인 신뢰 신호로 공식화하고 실증 검증하는 것이다. 본 연구는 명시적 승인 관계가 이론적으로 해석 가능하고 GMX V2 무기한 거래 성공에 정렬되며 거래 이력 기반 PageRank 기준선보다 계산적으로 더 확장 가능한 평판 지표를 산출할 수 있는지 보여 새로운 지식에 기여할 것으로 위치한다:


\proposalSubheading{1. 이론적 혁신}

\begin{itemize}
\item[\textbullet] \textbf{보증 기반 그래프 모델링:} ERC-20 approve 및 permit 이벤트에서 파생된 방향성 가중 보증 그래프를 도입함으로써, EndorseRank는 원래 PageRank 인용 패러다임을 명시적 온체인 신뢰 신호 포착으로 확장한다. 토큰 송금에서 승인 링크로의 이 개념적 전환은 블록체인 맥락에서 네트워크 중심성의 이론적 토대를 넓힌다.
\item[\textbullet] \textbf{시간적 단순화:} 정교한 시간 감쇠 함수를 요구하는 기존 접근과 달리, EndorseRank는 ERC-20 승인의 덮어쓰기 의미론을 활용하여 시간 기반 가중 없이 지속적 신뢰 관계를 모델링하는 더 깔끔하고 원칙적인 틀을 제공한다.
\item[\textbullet] \textbf{레버리지 노출 검증 틀:} 송금 in-degree/in-value만이 아닌 GMX V2 청산 성공 횟수 및 실현 이익 대리변수에 대해 평판 점수를 검증함으로써, 그래프 기반 평판을 Arbitrum One의 담보화된 무기한 거래 결과에 연결한다.
\end{itemize}
\proposalSubheading{2. 방법론적 기여}

\begin{itemize}
\item[\textbullet] \textbf{효율적 계산:} EndorseRank는 더 희소한 승인 그래프가 송금 기반 PageRank와 동등하거나 우수한 거래 성공 정렬의 평판 점수를 산출하면서 계산 시간, 메모리 사용, 수렴 반복을 줄일 수 있음을 보여준다. 이 효율성 향상은 자원 제약 DeFi 플랫폼에서 대규모 거의 실시간 평판 점수 산출을 더 실현 가능하게 한다.
\item[\textbullet] \textbf{재현 가능한 프레임워크:} Arbitrum One 데이터 추출부터 PageRank 구현 및 순위 상관 평가까지 완전히 명시된 오픈소스 파이프라인을 제공하여, 다른 연구자와 실무자가 다양한 프로토콜 및 생태계에서 EndorseRank 방법론을 검증, 확장, 적용할 수 있게 한다.
\end{itemize}
\proposalSubheading{3. 실무적 영향}

\begin{itemize}
\item[\textbullet] \textbf{리스크 관리 개선:} GMX 무기한 거래 성공과 정렬되는 온체인 보증 구조를 갖는 지갑을 더 정확히 식별함으로써, EndorseRank는 담보화 및 담보 부족 상품의 차별화된 리스크 파라미터를 지원할 수 있다. 이는 DeFi에서 상당한 유동성을 해제하고 금융 서비스 접근을 넓힐 잠재력이 있다.
\item[\textbullet] \textbf{프로토콜 통합:} EndorseRank의 경량 특성은 온체인 또는 오프체인 리스크 오라클, 거버넌스 대시보드, 규정 준수 도구에 통합하기에 적합하여 DeFi 프로토콜, 무기한 거래소, 수탁 서비스, 규제 기관에 실행 가능한 정보를 제공한다.
\end{itemize}
\proposalSubheading{4. 더 넓은 함의}

\begin{itemize}
\item[\textbullet] \textbf{교차 프로토콜 적용 가능성:} 본 연구는 Arbitrum One과 GMX V2에 집중하지만, EndorseRank 접근은 모든 ERC-20 호환 네트워크 또는 사용자 정의 토큰 표준에 일반화되어 신흥 L2 솔루션 및 대안 스마트 계약 플랫폼에 관련된다.
\item[\textbullet] \textbf{향후 연구의 토대:} 승인 네트워크와 GMX 파생 거래 성공 대리변수의 가치를 강조함으로써, 거버넌스 투표, 스테이킹 승인, 크로스체인 브리지 등 다른 온체인 신호의 추가 탐구와 복합 평판·신원 프레임워크 통합을 촉진한다.
\end{itemize}
